\textbf{1. Kiến thức cần nhớ}
\begin{itemize}
\item \textbf{HÌNH THOI}:
\begin{itemize}
\item \textbf{Định nghĩa}: Hình thoi là tứ giác có bốn cạnh bằng nhau ( = BC = CD = DA$). Hình thoi cũng là một hình bình hành.
\item \textbf{Tính chất}: Hình thoi có đầy đủ tính chất của hình bình hành và:
\begin{itemize}
\item Hai đường chéo \textbf{vuông góc} với nhau: $\boxed{AC \perp BD}$.
\item Hai đường chéo là các \textbf{đường phân giác} của các góc của hình thoi.
\end{itemize}
\item \textbf{Dấu hiệu nhận biết hình thoi}:
\begin{enumerate}
\item Tứ giác có bốn cạnh bằng nhau là hình thoi.
\item Hình bình hành có hai cạnh kề bằng nhau là hình thoi.
\item Hình bình hành có hai đường chéo vuông góc với nhau là hình thoi.
\item Hình bình hành có một đường chéo là đường phân giác của một góc là hình thoi.
\end{enumerate}
\end{itemize}

\item \textbf{HÌNH VUÔNG}:
\begin{itemize}
\item \textbf{Định nghĩa}: Hình vuông là tứ giác có bốn góc vuông và bốn cạnh bằng nhau.
\begin{itemize}
\item Hình vuông vừa là hình chữ nhật, vừa là hình thoi.
\end{itemize}
\item \textbf{Tính chất}: Hình vuông có tất cả các tính chất của hình chữ nhật và hình thoi:
\begin{itemize}
\item Bốn cạnh bằng nhau, các cạnh đối song song.
\item Bốn góc vuông bằng ^\circ$.
\item Hai đường chéo \textbf{bằng nhau, vuông góc với nhau tại trung điểm của mỗi đường} và là các \textbf{đường phân giác} của các góc:
80669\boxed{AC = BD,\ AC \perp BD,\ OA = OB = OC = OD}80669
\end{itemize}
\item \textbf{Dấu hiệu nhận biết hình vuông}:
\begin{enumerate}
\item Hình chữ nhật có hai cạnh kề bằng nhau là hình vuông.
\item Hình chữ nhật có hai đường chéo vuông góc với nhau là hình vuông.
\item Hình chữ nhật có một đường chéo là đường phân giác của một góc là hình vuông.
\item Hình thoi có một góc vuông là hình vuông.
\item Hình thoi có hai đường chéo bằng nhau là hình vuông.
\end{enumerate}
\end{itemize}
\end{itemize}

\textbf{2. Phương pháp giải các dạng}

%% DANG 8H3?14-1
\textbf{Dạng 1. Tính chất, dấu hiệu nhận biết hình thoi}
\begin{enumerate}
\item Sử dụng tính chất hai đường chéo vuông góc và cắt nhau tại trung điểm mỗi đường để tính độ dài theo định lí Pythagore:
80669\left(\dfrac{AC}{2}\right)^2 + \left(\dfrac{BD}{2}\right)^2 = a^280669
\item Diện tích hình thoi bằng nửa tích hai đường chéo: $\boxed{S = \dfrac{1}{2}d_1 d_2}$.
\item Chứng minh hình thoi: Thường chứng minh tứ giác là hình bình hành trước, sau đó chỉ ra hai cạnh kề bằng nhau hoặc hai đường chéo vuông góc.
\end{enumerate}

%% DANG 8H3?14-2
\textbf{Dạng 2. Tính chất, dấu hiệu nhận biết hình vuông}
\begin{enumerate}
\item Để chứng minh một tứ giác là hình vuông, con đường phổ biến nhất:
\begin{itemize}
\item Chứng minh là \textbf{hình chữ nhật} $\rightarrow$ chỉ ra thêm $ cạnh kề bằng nhau hoặc $ đường chéo vuông góc.
\item Hoặc chứng minh là \textbf{hình thoi} $\rightarrow$ chỉ ra thêm $ góc vuông hoặc $ đường chéo bằng nhau.
\end{itemize}
\item Đường chéo hình vuông cạnh $ có độ dài: $\boxed{d = a\sqrt{2}}$.
\end{enumerate}

%% DANG 8H3?14-3
\textbf{Dạng 3. Bài toán tổng hợp về tứ giác đặc biệt}
\begin{enumerate}
\item Sơ đồ chuyển hóa giữa các tứ giác:
80669\text{Tứ giác} \xrightarrow{\text{các cạnh đối } \parallel} \text{HBH} \xrightarrow{\text{1 góc vuông}} \text{HCN} \xrightarrow{2 \text{ cạnh kề } =} \text{Hình vuông}80669
80669\text{Tứ giác} \xrightarrow{\text{các cạnh đối } \parallel} \text{HBH} \xrightarrow{2 \text{ cạnh kề } =} \text{Hình thoi} \xrightarrow{\text{1 góc vuông}} \text{Hình vuông}80669
\item Dạng toán tìm điều kiện của tam giác ban đầu để tứ giác là hình chữ nhật, hình thoi, hình vuông:
\begin{itemize}
\item Giả sử tứ giác đã là hình chữ nhật (hoặc hình thoi, hình vuông), truy ngược lại tính chất cần có của tam giác ban đầu (ví dụ: tam giác phải cân, vuông, hoặc vuông cân).
\item Kết luận điều kiện và thử lại tính đúng đắn.
\end{itemize}
\end{enumerate}
